\begin{problem}[33届物理竞赛复赛]{60}
    一种拉伸传感器的示意图如图 a 所示：它由一半径为 $r_{2}$ 的圆柱形塑料棒和在上面紧密缠绕 N （N >> 1）圈的一层细绳组成；绳柔软绝缘，半径为 $r_{1}$ ，外表面均匀涂有厚度为 t （ $t << r_{1} << r_{2}$ ）、电阻率为 $\rho$ 的石墨烯材料；传感器两端加有环形电极（与绳保持良好接触）。未拉伸时，缠绕的绳可视为 N 个椭圆环挨在一起放置；该椭圆环面与圆柱形塑料棒的横截面之间的夹角为 $\theta$ （见图a），相邻两圈绳之间的接触电阻为 $R_{c}$ 。现将整个传感器沿塑料棒轴向朝两端拉伸，绳间出现 $n$ 个缝隙，每个缝隙中刚好有一整圈绳，这圈绳被自动调节成由一个未封闭圆环和两段短直线段（与塑料棒轴线平行）串接而成（见图b）。假设拉伸前后 $\theta$ 、 $r_1$ 、 $r_2$ 、 $\rho$ 、 $t$ 均不变。

    \begin{subproblem}
        求拉伸后传感器的伸长率 $\varepsilon$ （ $\varepsilon$ 是传感器两电极之间距离的伸长与其原长之比）和两环形电极间电阻的变化率；
    \end{subproblem}
    \begin{subproblem}
        在传感器两环形电极间通入大小为 $I$ 的电流，求此传感器在未拉伸及拉伸后，在塑料棒轴线上离塑料棒中点O距离为 $D$ （ $D$ 远大于传感器长度）的P点（图中未画出）处沿轴向的磁感应强度。
    \end{subproblem}

    已知：长半轴和短半轴的长度分别为 $a$ 和 $b$ 的椭圆的周长为 $\pi (3\frac{a + b}{2} -\sqrt{ab})$ ，其中 $b\neq 0$ 。
\end{problem}